\section{期望、线性性与指示变量}
\label{sec:05-Probability/04-Expectation-and-Moments/01}

三十个人的名单摆在面前，问一句：这份名单里平均有几对人同一天生日？

按部就班的做法是先求出``碰撞对数''这个随机变量的分布，再对取值加权求和。往下推两步就会发现路不好走——两对碰撞可能共用一个人，也可能毫不相干，情形随人数迅速分叉，光是写出通项就够累。

换个记法。三十个人两两配对，共 \(\binom{30}{2}=435\) 对。给第 \(k\) 对配一个只取 0 和 1 的变量 \(I_k\)：这一对同天生日记 1，否则记 0。碰撞总数于是写成 \(X=I_1+\cdots+I_{435}\)。假定生日在 365 天上均匀分布、各人相互独立，单独盯住任意一对，两人撞上同一天的概率是 \(1/365\)。所以

\[
\E[X]=\sum_{k=1}^{435}\E[I_k]=\frac{435}{365}\approx1.19.
\]

要留意的不是这个数字，而是这 435 个变量之间的关系有多脏。\(I_{AB}=1\) 与 \(I_{BC}=1\) 同时成立，就强制 \(I_{AC}=1\)；它们彼此高度纠缠，联合分布复杂得没人愿意写。上面那行计算对此一无所知，也不需要知道。

期望对加法的兼容不附带任何独立性条件。这是本节唯一要立住的判断，也是概率论里性价比最高的一条工具。

\subsection{定义只需要一个附加条件}

离散随机变量的期望就是按概率加权的平均：

\[
\E[X]=\sum_x x\,p_X(x),
\]

连续情形把求和换成积分，\(\E[X]=\int x f_X(x)\,dx\)。两种写法都附带同一个前提，\(\sum_x|x|p_X(x)<\infty\) 或 \(\int|x|f_X(x)\,dx<\infty\)，即绝对收敛。

这个前提不是形式上的洁癖。条件收敛的级数重排之后和会变，而 PMF 里各个取值本来就没有天然顺序；若不要求绝对收敛，期望值就会取决于你把哪一项写在前面。绝对收敛把顺序这个自由度锁死。

顺带澄清一个常见误解：期望不必是可能出现的取值。公平骰子的六个面是 \(1\) 到 \(6\)，期望是 \(3.5\)，这个点一次也掷不出来。它是分布的平衡点，不是最可能的结果，也不是``典型值''。

若关心的是 \(g(X)\) 而非 \(X\)，不必先去构造 \(g(X)\) 的分布：既然 \(X=x\) 发生时 \(g(X)\) 就等于 \(g(x)\)，直接按 \(X\) 的概率对变换后的数值加权即可，\(\E[g(X)]=\sum_x g(x)p_X(x)\)。这条记账方式有个夸张的名字，无意识统计学家定律（LOTUS），内容却只有上面这一句。

\subsection{线性性为什么与独立无关}

只要各项期望都存在且有限，

\[
\E\!\left[\sum_{i=1}^n a_iX_i+c\right]=\sum_{i=1}^n a_i\E[X_i]+c.
\]

两个变量的情形足以看清机制。把联合分布摊开，

\[
\begin{aligned}
\E[aX+bY]&=\sum_{x,y}(ax+by)\,P(X=x,Y=y)\\
&=a\sum_x x\underbrace{\sum_y P(X=x,Y=y)}_{=\,P(X=x)}
+b\sum_y y\underbrace{\sum_x P(X=x,Y=y)}_{=\,P(Y=y)}.
\end{aligned}
\]

用到的只有一件事：把联合分布沿一个方向求和，得到另一个变量的边缘分布。这是概率归一化的直接后果，与 \(P(X=x,Y=y)\) 能否分解成两个因子毫无关系。独立性在整段推演里从未登场，自然也就无从要求。

对比之下，\(\E[XY]=\E[X]\E[Y]\) 才是需要独立性（或至少不相关）撑着的式子。加法免费，乘法收费——这条界线值得记牢，下一节讨论方差时它会立刻发作。

\begin{center}
\begin{tikzpicture}[x=1cm,y=1cm,font=\small,>=stealth]
  \foreach \i/\v in {0/1,1/0,2/1,3/1,4/0}{
    \draw[black!55,fill=blue!12] (1.5*\i,0) rectangle (1.5*\i+1.0,1.0);
    \node at (1.5*\i+0.5,0.5) {\(\v\)};
  }
  \foreach \i/\k in {0/1,1/2,2/3,3/4,4/5}{
    \node[below] at (1.5*\i+0.5,-0.02) {\(I_{\k}\)};
  }
  \draw[dashed,black!55] (0.5,1.02) .. controls (1.25,1.95) .. (2.0,1.02);
  \draw[dashed,black!55] (2.0,1.02) .. controls (3.5,2.30) .. (5.0,1.02);
  \draw[dashed,black!55] (3.5,1.02) .. controls (4.25,1.95) .. (5.0,1.02);
  \node[black!60,font=\footnotesize] at (3.5,2.55) {虚线标出彼此不独立的配对};
  \node at (3.25,-0.62) {\(X=I_1+I_2+I_3+I_4+I_5\)};
  \draw[->,black!70] (3.25,-0.92) -- (3.25,-1.32);
  \node at (3.25,-1.62) {\(\E[X]=P(A_1)+P(A_2)+P(A_3)+P(A_4)+P(A_5)\)};
\end{tikzpicture}

\emph{一堆 0--1 变量相加，取期望时逐项落到各自事件的概率上。虚线画的依赖关系一根也没有断，它们只是不参与这笔账——求和与取期望这两个动作可以交换顺序，交换的理由与变量之间的关系无关。}
\end{center}

\subsection{把计数拆成指示变量}

对任意事件 \(A\)，指示变量 \(\ind_A\) 在 \(A\) 发生时取 1、否则取 0，它的期望干脆利落：

\[
\E[\ind_A]=1\cdot P(A)+0\cdot P(A^c)=P(A).
\]

与线性性拼在一起，就得到一个可以反复复用的动作。想求某个计数量的期望时，先问它数的是什么；把被数的对象一个个列出来，每个对象配一个指示变量；\(X\) 于是等于这些指示变量之和，\(\E[X]\) 等于一串概率相加。整套流程里唯一需要动脑的是单个对象的概率，而单个对象通常简单到可以口算。

开头的生日问题就是这个动作的第一例。再看两个。

\begin{example}[随机排列的不动点]
把 \(1,\ldots,n\) 随机打乱，问平均有几个元素留在原位。直接求这个数的分布要走错排公式和容斥原理，篇幅可观。改成指示变量：令 \(I_i\) 表示元素 \(i\) 是否落回位置 \(i\)。均匀随机排列下每个位置地位相同，\(\E[I_i]=1/n\)，于是 \(\E[X]=n\cdot(1/n)=1\)。无论 \(n\) 是 10 还是 10000，平均都恰好有一个元素没动。
\end{example}

\begin{example}[不放回抽牌]
52 张牌里不放回抽 5 张，问平均抽到几张 A。走分布的话要写出 Hypergeometric 的 PMF（见\hyperref[sec:05-Probability/06-Common-Distributions/02]{``不放回抽样与 Hypergeometric 分布''}）再逐项求和，四五项、每项三个组合数。走指示变量：第 \(i\) 张是不是 A，记作 \(I_i\)。不放回让这五个变量互相牵制，但每个位置单独看都是从 52 张里随机取一张，\(\E[I_i]=4/52\)，所以 \(\E[X]=5\times\frac1{13}=\frac5{13}\)。组合数一个也没出现。
\end{example}

拆分的零件不一定得是指示变量，任何能把目标写成和的分解都吃这条性质。优惠券收集问题是常被引用的例子：\(n\) 种赠品，每次随机得到一种，集齐全部平均要抽多少次？把总次数 \(T\) 按``已集齐 \(k-1\) 种''分段，第 \(k\) 段的等待次数 \(T_k\) 服从成功概率 \((n-k+1)/n\) 的几何分布，期望是 \(n/(n-k+1)\)。相加得

\[
\E[T]=\sum_{k=1}^{n}\frac{n}{n-k+1}=n\Bigl(1+\frac12+\cdots+\frac1n\Bigr)=nH_n\approx n\ln n.
\]

\(n=50\) 时约 225 次。这里各段等待恰好是独立的，但独立与否并不影响相加这一步的合法性，只影响后面若要算方差时会不会多出交叉项。

在数据工作里，这个动作最常见的形态是批量指标。一天的日志里有多少条请求超时、多少个会话触发了推荐、多少对用户被判为重复账号——这些都是``对每个对象查一次，再数有多少个 1''。样本之间往往共享用户、共享时段、共享模型版本，谈不上独立；但只要你要的是期望总量，逐条把单条概率加起来就够了，不需要为相关性做任何修正。要估计的东西一旦从``平均多少''变成``会不会超过阈值''，情况立刻改变，那属于集中不等式的地盘（见\hyperref[ch:043]{``变换与概率界：不知道完整分布时还能说什么''}）。

\subsection{期望也有失效的时候}

不是每个分布都有期望。取 \(P(X=2^k)=2^{-k}\)，\(k=1,2,\ldots\)，概率之和是 1，是个合法分布，但

\[
\E[X]=\sum_{k\ge1}2^k\cdot2^{-k}=\sum_{k\ge1}1=\infty.
\]

更麻烦的是正负两侧同时发散的情形，此时 \(\E[X]\) 连``等于无穷''都算不上，直接没有定义——\(\infty-\infty\) 不是数值。现实里的重尾现象（保险单笔赔付、社交网络的转发量、部分金融资产的收益）确实可能落在这一档，样本均值会随着数据增多不断漂移而不收敛，参见\hyperref[sec:05-Probability/06-Common-Distributions/07]{``重尾分布：当均值与方差失效时''}。写得出分布不等于矩存在，这是用任何一阶量之前该先核一遍的事。

即便期望存在且有限，它也只报告中心位置。两个方案的期望收益都是 10，一个是稳稳的 10，另一个以 \(99.9\%\) 的概率给 0、以 \(0.1\%\) 的概率给 10000，从期望看不出任何区别。要把它们分开，得先有一个度量``散得有多开''的量。

拆成一串再逐项相加，这条路对期望永远畅通。下一节把同样的问题问给方差：\(\Var(X+Y)\) 是否等于 \(\Var(X)+\Var(Y)\)？答案在最简单的例子上就会崩，而崩掉的那部分恰好定义了本单元后半段的主角。
